% !TeX root = ../深度学习笔记.tex
% 由主文档通过 \input 载入；请编译主文档。
\setcounter{chapter}{23} % 保留原书章号。
\chapter{效率、部署与综合项目}
\label{chap:dl-projects}
模型的实际价值同时取决于任务表现、计算资源和输入条件。一个可复核项目应让读者知道结果怎样得到、代价是什么、在哪些情况下会失败，而不只保留最好的一次指标。

\section{计算量、内存与性能测量}
设模型有 $p$ 个参数，每个数值用 $b$ 字节，则仅权重占 $bp$ 字节；训练还要保存梯度、优化器状态、中间激活和临时工作区。以全部使用四字节数值的 Adam 为例，权重、梯度、一阶矩和二阶矩共约 $16p$ 字节，未计激活、主权重副本或实现开销。长序列推断还需考虑 KV cache，不能只根据参数文件大小选择设备。

吞吐量是单位时间处理的样本或词元数，延迟是一次请求从输入到输出的时间；增加批量可提高吞吐，也可能增加单请求等待。测量应注明硬件、数值精度、批量、输入长度、预热次数和同步方式，并区分纯模型与包含预处理、传输、后处理的端到端时间。理论乘加量只能解释部分成本，不能代替实测。

\section{混合精度、累积梯度与并行训练}
\emph{混合精度（mixed precision）}在适合的运算中使用较低精度，同时对敏感计算或状态保留较高精度。损失缩放将损失乘 $S>0$，反向后再除以 $S$，在精确算术下不改变梯度，可缓解某些格式的下溢；溢出检测和缩放调整属于实际算法的一部分。激活重计算在反向时重做部分前向，交换计算与内存，要求重算时随机性与状态行为保持一致。

若总共 $N=\sum_jB_j$ 个样本被分为微批量，且参数在整个累积期间固定、模型没有批量相关操作，则全体平均梯度等于 $\sum_j(B_j/N)\nabla J_j$，其中 $J_j$ 为第 $j$ 个微批量平均损失。应清零一次、按这些权重累积、最后更新一次。批量归一化使不同分组的目标不同，逐微批量更新又改变了求导点，都不能直接视为同一次大批量更新。

数据并行在多个设备复制模型、分配数据并聚合梯度。各设备样本数相同可平均设备梯度，不相同时按样本数加权；同步还增加通信开销。比较单设备和多设备实验时，说明全局批量、更新次数及学习率是否改变，避免把训练协议变化完全归因于硬件加速。

\section{模型压缩与推断部署}
\emph{量化（quantization）}把数值映到较少的离散水平，例如 $q=\operatorname{clip}(\operatorname{round}(x/s)+z_0,q_{\min},q_{\max})$、$\widehat x=s(q-z_0)$，其中 $s>0$ 为尺度、$z_0$ 为零点。舍入与截断引入误差，范围估计和敏感层处理影响表现。\emph{剪枝（pruning）}移除权重或结构单元，稀疏权重只有得到相应算子支持才可能加速。

\emph{知识蒸馏（knowledge distillation）}用教师分布训练学生。设温度 $\tau>0$ 下的教师与学生输出为 $p_T^{(\tau)},p_S^{(\tau)}$，可取 $\lambda\ell_{\mathrm{label}}+(1-\lambda)\tau^2D_{\mathrm{KL}}(p_T^{(\tau)}\|p_S^{(\tau)})$，$0\leq\lambda\leq1$；温度和缩放约定属于目标定义。教师错误也可能被传递，必须保留独立评价。部署时固定预处理、推断模式和后处理，验证压缩前后同一输入的输出及端到端成本。

\section{三个基础项目与一个进阶项目}
\textbf{图像分类项目。}选择带明确来源和划分规则的小型数据，先实现多数类与线性基线，再比较 MLP、CNN 和一个残差变体。固定主要训练预算，分别改变网络结构和增强，不同时混改；报告分类指标、参数量、训练时间及典型错误。先验证一个小批量的形状与梯度，再运行完整训练。

\textbf{序列项目。}构造延迟复制、括号类型预测等可自动检查标签的任务，或使用固定来源的短文本。比较 RNN 与小型 Transformer，核查目标移位和未来信息屏蔽；在训练长度内外分别评价，报告任务正确率或困惑度，并单列生成延迟。样本生成规则与最大长度作为实验条件公开。

\textbf{迁移项目。}固定源模型，在多个目标标注预算下比较从头训练、线性探测、部分与全量微调。所有方案使用一致测试集和相近选择预算，报告预训练来源、目标性能、源能力保留及成本。\textbf{进阶生成项目。}在已知模式的二维分布上任选 VAE、GAN、流或扩散，先核对核心概率或密度公式，再比较采样与真实分布的模式比例、局部质量和成本；没有训练结果前不预设某模型胜出。

\section{项目验收、报告与后续学习}
项目报告应把问题定义、数据与划分、算法目标、配置、基线、对照、主要指标和成本连接成可复核的论证。清楚标注理论保证与经验发现，说明随机种子、检查点选择和评价脚本的规则；保留失败和无提升结果，而不是只汇报最优曲线。一个完整实验至少要让别人能够确定输入、复现主要计算并验证评价含义。

\begin{example}[训练状态的存储估计]
某模型有一千万个参数，以四字节数值保存参数、梯度和 Adam 两个矩状态。估算这些张量的总存储。
\end{example}
\begin{solve}
四组张量共 $4\times4\times10^7=1.6\times10^8$ 字节，约为十进制 $160$ MB。该数不是实际峰值内存：激活、临时工作区、数据批量和可能的参数副本尚未计入。是否能运行仍需结合输入形状与实现测量。
\end{solve}
\begin{enumerate}
\item 三个微批量大小为 $4,4,2$，写出正确的累计损失权重，说明为何全部除以三一般不对。
\item 固定输入和设备，测量两种批量大小的吞吐量与延迟，解释二者为何可能给出不同排序。
\item 选择上述一个项目，先写出可反驳的假设、匹配条件和判定标准，再实施实验，最后逐条说明证据支持了什么。
\end{enumerate}

